\section{The nonlinear map and the material isomorphism}\label{inf:sec-material}

The chart in this section uses a harmonic extension on the whole ball.
It is different from the homogeneous radial stretching used to derive
the spectral flow. Both charts describe the same boundary domain at a
fixed centre; only its spectrum is passed from one description to the
other.

\subsection{A smooth scalar chart on the ball}

For a normalized boundary shape $f\in\Ec_1$ put
\begin{equation}\label{inf:scalar-chart}
 H_f=\HP f,\qquad q=1+H_f,\qquad d=q+\Euler q,
 \qquad \Phi_f(y)=q(y)y.
\end{equation}
For a field $U$ on the reference ball define
\begin{equation}\label{inf:scalar-pullback-definition}
 P_{\Phi_f}\Delta U=
       [\Delta(U\circ\Phi_f^{-1})]\circ\Phi_f,
 \qquad B_{\Phi_f}=P_{\Phi_f}\Delta+R^2.
\end{equation}
For a smooth real chart, $B_{\Phi,D}$ denotes the restriction of
$B_\Phi$ to the invariant Dirichlet graph domain
$H^2(B)\cap H^1_0(B)$.
Fix the analytic chart radius $\varepsilon_0=1/2$ and the
working radius $\varepsilon_{\mathrm{work}}=1/4$.
For every real $f\in\Ec_1$ with $\nrm f_{\Ec_1}<1/2$, the map
$\Phi_f$ is a diffeomorphism of the closed ball onto the region
with radial boundary $1+f$.
Indeed $|q-1|\le\nrm f_{\Ec_1}$ and
$|d-1|\le2\nrm f_{\Ec_1}<1$ on the ball, so $q,d>0$.
On each ray, $\Phi_f(\zeta\omega)=\zeta q(\zeta\omega)\omega$
keeps the direction fixed and has strictly positive radial
derivative $d$. Distinct rays remain distinct and the radial
map is one-to-one on each ray. Its Jacobian determinant is
$q^{2p-1}d>0$, including at the origin. The inverse function
theorem and the radial endpoint give the global
identification. By Lemma~\ref{inf:space-embedding}, the harmonic
extension and chart are real analytic on a neighbourhood of
the closed ball. Both axes and the origin are ordinary points
of this chart.

\begin{lemma}[Exact scalar pullback]\label{inf:pullback}
Let $w=\nabla q$, $M_q=w\cdot w$ and
\begin{equation}\label{inf:pullback-metric}
 A_q=q^{-2}\left\{I-\frac{y\otimes w+w\otimes y}{d}
                         +\frac{M_q}{d^2}y\otimes y\right\},
 \qquad C_q=A_q:D^2q.
\end{equation}
Then
\begin{align}
 P_{\Phi_f}\Delta U={}&q^{-2}\Delta U
 -\frac2{q^2d}\nabla q\cdot\nabla\Euler U
 +\frac{M_q}{q^2d^2}\Euler(\Euler+1)U
 \label{inf:scalar-pullback}\\
 &+\left(\frac{2M_q}{qd^3}-\frac{C_q}{d}\right)\Euler U,\nonumber\\
 C_q={}&q^{-2}\left\{-\frac{2(\nabla q\cdot\nabla\Euler q-M_q)}d
                +\frac{M_q\Euler(\Euler-1)q}{d^2}\right\}.
 \label{inf:Cq-identity}
\end{align}
The second identity uses harmonicity of $q$.
\end{lemma}
\begin{proof}
Since $D\Phi_f=qI+y\otimes w$, its inverse is
$q^{-1}(I-y\otimes w/d)$ and
$D\Phi_f^{-1}D\Phi_f^{-T}=A_q$.
The coordinate change formula for a scalar Laplacian is
\[
 P_{\Phi_f}\Delta U=A_q:D^2U
   -D\Phi_f^{-1}(A_q:D^2\Phi_f)\cdot\nabla U.
\]
Now $A_q:D^2\Phi_f=yC_q+2A_qw$, and multiplication by the
inverse matrix gives
\[
 D\Phi_f^{-1}(A_q:D^2\Phi_f)
 =\frac{2w}{q^2d}
   +y\left\{\frac{C_q}d-\frac{2M_q}{q^2d^2}
                             -\frac{2M_q}{qd^3}\right\}.
\]
On the other hand,
\[
 A_q:D^2U=q^{-2}\left\{\Delta U-\frac2d
 (\nabla q\cdot\nabla\Euler U-\nabla q\cdot\nabla U)
             +\frac{M_q}{d^2}\Euler(\Euler-1)U\right\}.
\]
The terms with $\nabla q\cdot\nabla U$ cancel, giving
\eqref{inf:scalar-pullback}.
Apply the last formula for $A_q:D^2U$ to $U=q$ and use
$\Delta q=0$ to get \eqref{inf:Cq-identity}.
All computations are identities of smooth functions; extension to
the coefficient spaces is established next.
\end{proof}

\begin{proposition}[Analytic map with a polynomial second derivative]
\label{inf:nonlinear-map}
For $g\in\Gc_p$ and $\nrm f_{\Ec_1}<\varepsilon_0$, define
\begin{equation}\label{inf:nonlinear-definition}
 \Fnl(g,f)=B_{\Phi_f}(1+Kg).
\end{equation}
This is an analytic map
$\Gc_p\times\{f\in\Ec_1:\nrm f_{\Ec_1}<\varepsilon_0\}\to\Cc$,
affine in $g$,
and on $\nrm f_{\Ec_1}\le\varepsilon_{\mathrm{work}}=1/4$ it satisfies
\begin{equation}\label{inf:nonlinear-second}
 \nrm{D^2\Fnl(g,f)}\le Cp(1+\nrm g_\Cc).
\end{equation}
For real $f$, its coefficient extension is exactly the weak scalar
pullback of $1+Kg\in H^2(B)$. Its complexification is the analytic
extension of the same algebraic operator formulas.
\end{proposition}
\begin{proof}
The harmonic module estimate gives
$\nrm{M_{q-1}}\le\nrm f_{\Ec_1}$. For $d-1=H_f+\Euler H_f$,
the same estimate applied to each solid coefficient gives
\[
 \nrm{M_{d-1}}
 \le\sum_{j\ge0}(1+2j)2^j\wt(2j)|f_j|
 \le2\sum_{j\ge0}(1+j)2^j\wt(2j)|f_j|
 =2\nrm f_{\Ec_1}.
\]
Their inverse powers are norm-convergent Neumann series on
$\nrm f_{\Ec_1}<1/2$. On the closed working ball of radius $1/4$,
differentiating these series any fixed number of times gives
uniformly bounded multilinear multiplier maps.
The same series converge uniformly on the closed ball, by
\eqref{inf:shape-C2}, to the indicated pointwise inverse powers.

Apply \eqref{inf:scalar-pullback} to $U=1+Kg$.
All derivatives of the constant vanish. The term $q^{-2}\Delta Kg$
is $q^{-2}g$.
Lemma~\ref{inf:gradient-profiles} controls
$\nabla q\cdot\nabla\Euler K$ with one shape moment.
The operator $\Euler(\Euler+1)K$ is bounded by
\eqref{inf:poisson-operator-bounds}, and its multiplier $M_q$
is bounded by \eqref{inf:gradient-multiplier} with norm at most
$Cp\nrm f_{\Ec_1}^2$.
For the last term group
\[
 C_q\Euler K=q^{-2}\left\{
 -2d^{-1}(\nabla q\cdot\nabla\Euler q-M_q)\Euler K
 +M_qd^{-2}[\Euler(\Euler-1)q]\Euler K\right\}.
\]
The first grouped product is bounded by
\eqref{inf:mixed-shape-clamped} and
\eqref{inf:gradient-multiplier}; the last one by
\eqref{inf:second-shape-clamped} followed by the bounded multiplier
$M_q$. The scalar multipliers commute, so the high shape derivative
is always placed immediately next to $\Euler K$ on its clamped
source, as in Proposition~\ref{inf:grouped-shapes}.
Every numerator is thus a bounded linear, bilinear, or trilinear
map from shapes to operators $\Gc_p\to\Cc$, of norm at most $Cp$.
The inverse multiplier series prove analyticity of the operator
coefficient $L_f\in\mathcal B(\Gc_p,\Cc)$ in
$\Fnl(g,f)=R^2+L_fg$.
Its first two shape derivatives have norm at most $Cp$ on the
closed working ball of radius $1/4$. The two mixed derivatives and the pure shape second
derivative give \eqref{inf:nonlinear-second}; the pure source
second derivative is zero. The term $R^2(1+Kg)$ is shape independent
and contributes no second derivative.

To identify the extension, first use polynomial sources and finite
shape sums, where the identity is classical. Approximate any
$f\in\Ec_1$ by finite sums in its norm; the charts converge in
$C^2$ and have uniformly positive $q,d$ on any fixed smaller shape
ball. The explicit inverse matrix is therefore uniformly
bounded there. For each fixed $p$, their scalar
pullbacks converge as maps $H^2(B)\to L^2(B)$ by the coordinate
formula. Next approximate $g$ in $\Cc$, hence in $L^2$.
Dirichlet regularity and uniqueness give convergence of $Kg$ in
$H^2(B)$. The bounded coefficient maps converge in $\Cc$, hence
in $L^2$, to the same pullback. This proves the assertion on the
completed spaces.
\end{proof}

\subsection{Coefficient size of the known centre}

\begin{lemma}[Fixed radial modes in coefficient norm]
\label{inf:radial-centre-norm}
Fix the finite mode set of the order-$N$ centre in
Theorem~\ref{inf:centres}. For each fixed moment $k$ and
$p\ge P_{N,k}$, its regular radial modes,
rescaled to the unit ball, have $\Cc^{[k]}$ norms at most
$C_{N,k}p^{k+1}$. The bound is uniform in $r$ and
$0<|\theta|\le20$.
\end{lemma}
\begin{proof}
For one fixed $j$ put $\nu=p+2j-1$ and write its normalized mode as
$\zeta^{1-p}J_\nu(R\zeta)G_j/d_j^{\mathrm{mode}}$.
Here $d_j^{\mathrm{mode}}$ is either $J_\nu(R)$ for value normalization or
$J'_\nu(R)-(p-1)J_\nu(R)/R$ for physical normal normalization.
The finite recurrence \eqref{inf:finite-q-recurrence} shows that
$d_j^{\mathrm{mode}}/J_{p-1}(R)$ has a nonzero analytic limit in every case.
In particular the resonant seed is normalized by its normal trace.

In the solid basis the coefficients of this mode are
\begin{equation}\label{inf:Bessel-Jacobi-coefficients}
 c_s=\frac{2(\nu+1+2s)(-1)^sJ_{\nu+1+2s}(R)}{Rd_j^{\mathrm{mode}}}.
\end{equation}
To derive them, Rodrigues' formula and $s$ integrations by parts
in $t=\zeta^2$ give the integral identity
\[
 \int_0^1\zeta^{\nu+1}J_\nu(R\zeta)J_s^\nu(\zeta^2)\,d\zeta
 =\frac{(-1)^s}{R}J_{\nu+2s+1}(R).
\]
One may evaluate the final integral by the absolutely convergent
Bessel power series and the beta integral. The derivative identity
used in the integrations is
$\partial_t^s(t^{-\nu/2}J_\nu(R\sqrt t))
=(-R/2)^st^{-(\nu+s)/2}J_{\nu+s}(R\sqrt t)$.
Dividing by the radial norm in \eqref{inf:radial-norm} proves
\eqref{inf:Bessel-Jacobi-coefficients}.
Parseval and the Lommel integral at a general argument give
\begin{equation}\label{inf:centre-parseval}
 \sum_{s\ge0}\frac{|c_s|^2}{\nu+1+2s}
 =\frac{J_\nu(R)^2-J_{\nu-1}(R)J_{\nu+1}(R)}{(d_j^{\mathrm{mode}})^2}
 \le C_{j,\mathrm{rad}}.
\end{equation}
The general Lommel identity follows by differentiating
$x^2(J_\nu(x)^2-J_{\nu-1}(x)J_{\nu+1}(x))/2$ and using the
recurrences. The bound follows since every numerator value is a
fixed recurrence shift of $J_{p-1}(R)$ and the normalization
quotient has a nonzero limit.

For $s\le S=C_{j,\mathrm{rad}}'p+O_j(1)$, with a fixed $C_{j,\mathrm{rad}}'$ large enough that
$\nu+1+2S\ge2R$, Cauchy--Schwarz in
\eqref{inf:centre-parseval} gives
\[
 \sum_{s\le S}\wt(2j+2s)|c_s|\le C_{j,\mathrm{rad}}p.
\]
There are $O_j(p)$ terms; their denominator reciprocals in the
Cauchy--Schwarz factor are inverted to quantities $O_j(p)$, and
the weights are bounded by a constant by \eqref{inf:weight-properties}.
A fixed degree moment adds at most $C_{j,k}p^k$.

For the remaining tail use the positive ratio bound
$J_{m+1}(R)/J_m(R)<R/(m+1)$ when $m\ge2R$.
Here positivity follows from $j_{m,1}>m$.
For a proof of the ratio bound, its Riccati equation is
$r'=1-(2m+1)r/x+r^2$, with
$r(x)\sim x/(2m+2)$ at zero. At a first contact with
$x/(m+1)$ for $x\le m$, the difference of derivatives would be
$-1+x^2/(m+1)^2<0$, excluding upward contact.
Applying the bound twice to \eqref{inf:Bessel-Jacobi-coefficients}
gives $|c_{s+1}/c_s|<R^2/(m(m+1))\le1/4$ for
$m=\nu+1+2s\ge2R$.
The neighbouring grading ratio is at most $e^{2/p}$, and the
moment ratio tends uniformly to one as $p$ grows for fixed $k$.
Thus the weighted tail is geometric with ratio below $1/2$ for
large $p$, and is bounded by the head estimate. The angular factor
$2^j$ is bounded for the fixed list of modes.
\end{proof}

\begin{proposition}[The deformed centre and its source]
\label{inf:centre-source-norm}
Let $f_N=h_N/R$, $\Phi_N=\Phi_{f_N}$, and
\begin{equation}\label{inf:clamped-centre}
 U_*=u_N\circ(R\Phi_N),\qquad U_0=U_*-\ell_N,
 \qquad g_0=\Delta U_0,\qquad z_0=(g_0,f_N).
\end{equation}
For every fixed moment $k$ and $p\ge P_{N,k}$,
$g_0\in\Gc_p$, $U_0=1+Kg_0$, and
\begin{equation}\label{inf:centre-norm-bounds}
 \nrm{U_*}_{\Cc^{[k]}}+\nrm{U_0}_{\Cc^{[k]}}
 \le C_{N,k}p^{k+1},\qquad
 \nrm{g_0}_{\Cc^{[k]}}\le C_{N,k}p^{k+5}.
\end{equation}
For $p\ge P_{N,0}$, on the fixed product region
\begin{equation}\label{inf:nonlinear-working-region}
 \{g\in\Gc_p:\nrm{g-g_0}_{\Cc}\le1\}
 \ \times\ \{f\in\Ec_1:\nrm f_{\Ec_1}\le1/4\},
\end{equation}
\begin{equation}\label{inf:p6-second}
 \sup\nrm{D^2\Fnl}\le C_Np^6.
\end{equation}
\end{proposition}
\begin{proof}
For each known centre mode, fix $p$ and put $\nu=p+2j-1$.
Its radial factor in the solid basis is the entire function of
the complex squared-radius variable $t$,
\[
 F_{\nu,j}(t)=\frac1{d_j^{\mathrm{mode}}}
 \sum_{h\ge0}\frac{(-1)^h(R/2)^{\nu+2h}}{h!\,\Gamma(\nu+h+1)}t^h,
 \qquad t\in\C,
\]
where $\Gamma$ is Euler's gamma function
\cite[Eq.~10.2.2]{inf:dlmf}. For real $t>0$, this is
$t^{-\nu/2}J_\nu(R\sqrt t)/d_j^{\mathrm{mode}}$.
The radial expansion in Lemma~\ref{inf:radial-centre-norm} is
$\sum_{s\ge0}c_sJ_s^\nu(t)$ on $0\le t\le1$ in weighted $L^2$.
For $m=\nu+1+2s\ge2R$, the ratio bound proved in that lemma gives
$|c_{s+1}|\le R^2|c_s|/[m(m+1)]$. Since $m\ge2s$, iteration
from a fixed $s_0\ge1$ with $\nu+1+2s_0\ge2R$ gives
\[
 |c_s|\le |c_{s_0}|(R^2/4)^{s-s_0}
                 \left(\frac{(s_0-1)!}{(s-1)!}\right)^2,
 \qquad s\ge s_0.
\]
For every fixed $A>0$, Taylor's formula at one and
\eqref{inf:Jacobi-endpoint-derivatives} give
\[
 \sup_{|t-1|\le A}|J_s^\nu(t)|
 \le\sum_{h=0}^s\frac{[A s(2s+\nu+1)]^h}{(h!)^2}
 \le\exp\{2\sqrt{A s(2s+\nu+1)}\}.
\]
For fixed $p,j,A$, this envelope is at most exponential in $s$,
so the factorial bound makes the radial series locally uniformly
convergent on $\C$. Its sum is entire. On $[0,1]$ uniform
convergence and the original $L^2$ expansion identify that sum
with $F_{\nu,j}$; continuity turns the almost-everywhere equality
into equality on the interval. The identity theorem gives
equality on $\C$, including at the deformed radial arguments.

Write $q=1+\HP f_N$ and $e=q^2-1$.
The finite centre shape and the module estimate give
$\nrm{M_e}\le C_N\tau p^{-2}=:\epsilon_{\mathrm{cmp},p}$.
For one radial basis polynomial the exact Taylor identity is
\[
 J_s^\beta(q^2t)=\sum_{m=0}^s\frac{e^m}{m!}
                                      t^m\partial_t^mJ_s^\beta(t).
\]
The derivative columns and the positive $t$ stencils show that
$t^m\partial_t^mJ_s^\beta$, whose degree is no greater than $s$,
has coefficient sum
$(s)_{\mathrm{fall},m}(s+\beta+1)_m/m!$.
It follows that its relative composition norm is at most
\begin{equation}\label{inf:composition-bound}
 \sum_{m=0}^s\epsilon_{\mathrm{cmp},p}^m
 \frac{(s)_{\mathrm{fall},m}(s+\beta+1)_m}{(m!)^2}
 \le\exp\{2\sqrt{\epsilon_{\mathrm{cmp},p}s(2s+\beta+1)}\}.
\end{equation}
Indeed bound the product by $[s(2s+\beta+1)]^m$ and use
$\sum z^m/(m!)^2\le e^{2\sqrt z}$, which follows by comparing
with the even terms of the exponential using
$\binom{2m}m\le4^m$.
The extra solid factor is $q^{2j}$, a bounded multiplier because
$j$ belongs to the fixed finite set of centre modes.

On the $s=O_N(p)$ head the right side of
\eqref{inf:composition-bound} is bounded by $C_N$.
On the tail it is at most $\exp(C_N(s+p)/p)$.
Multiply the geometric Bessel tail in
Lemma~\ref{inf:radial-centre-norm} by this envelope and by the
weight ratio $e^{2/p}$. For large $p$ the resulting geometric
ratio is still below $1/2$. Thus each composed mode series
converges in $\Cc$ with norm at most $C_Np$.
The finite polynomial chart increases each total degree by at
most a fixed multiple depending on $N$. Therefore each output
moment is bounded by a constant times the corresponding input
moment; the same argument gives convergence and the bounds in
all fixed moments $k$.
For this fixed $p$, the radial entire-function identity makes
the composed partial sums converge uniformly on the closed ball
to $S_j(qy)F_{\nu,j}(q^2t)$: the arguments $q^2t$ range over a
compact set, and $S_j(qy)$ is a fixed polynomial. The polynomial
partial sums also converge in $\Cc$ by the preceding coefficient
estimate, hence in $L^2(B)$ by Lemma~\ref{inf:space-embedding}.
Uniform convergence on the closed ball gives the same $L^2$ limit.
Uniqueness of that limit and injectivity of $\Cc\hookrightarrow L^2(B)$
identify the coefficient sum with the composed mode.
Adding the finite list of centre modes proves the asserted
identification and bounds for $U_*$.
The finite field amplitudes are uniformly bounded by the centre
construction. Subtracting the lift uses \eqref{inf:lift-estimates}
and proves the assertion for $U_0$.
The Laplacian bound \eqref{inf:Delta-moment}, with moments included,
then gives $\nrm{g_0}_{\Cc^{[k]}}\le C_{N,k}p^{k+5}$.

By construction the two traces of $U_0$ are exactly $1,0$.
Dirichlet Poisson uniqueness yields $U_0=1+\KD g_0$, and
\eqref{inf:poisson-trace} gives $g_0\in\Gc_p$.
Finally \eqref{inf:nonlinear-second} and the $p^5$ source bound
give \eqref{inf:p6-second} on the product region
\eqref{inf:nonlinear-working-region}. Its exponent does not depend
on $N$.
\end{proof}

\subsection{Residuals in the same coefficient space}

\begin{proposition}[The residual and its material derivative]
\label{inf:residual-estimates}
At the centre $z_0$ of \eqref{inf:clamped-centre}, put
$E_0=\Fnl(z_0)$. For every fixed moment $k$ and $p\ge P_{N,k}$,
\begin{equation}\label{inf:residual-moments}
 \nrm{E_0}_{\Cc^{[k]}}\le C_{N,k}\theta^2p^{2-N},\qquad
 \nrm{\Euler E_0}_{\Cc^{[k]}}\le C_{N,k}\theta^2p^{3-N}.
\end{equation}
For a shape variation $\eta\in\Ec_1$, define
\begin{equation}\label{inf:residual-material}
 \mathcal V_{\mathrm{res}}(\delta g,\eta)
       =(\HP\eta)d_0^{-1}\Euler E_0,
 \qquad d_0=1+\HP f_N+\Euler\HP f_N.
\end{equation}
Then, for $p\ge P_N$,
\begin{equation}\label{inf:residual-material-bound}
 \nrm{\mathcal V_{\mathrm{res}}}_{\Xc_p\to\Cc}
 \le C_N\theta^2p^{3-N}.
\end{equation}
\end{proposition}
\begin{proof}
The interior equation of the finite centre is exact, so
$B_{\Phi_N}U_*=0$ and $E_0=-B_{\Phi_N}\ell_N$.
The lift occupies only radial layers zero and one.
On these layers the operators $\Delta$, $\Euler^2$ and
$\nabla a\cdot\nabla\Euler$, for any fixed solid polynomial $a$,
have coefficient bounds $C_a p^2$ times a fixed number of angular
moments.
On layer one,
\[
 \Euler\Psi_{j,1}=(2j+2)\Psi_{j,1}
                            +2(\beta_j+1)\Psi_{j,0}.
\]
On a solid monomial $S_l t^h$ of bounded radial power,
$\Delta(S_lt^h)=4h(h+2l+p-1)S_lt^{h-1}$ costs one factor $p+l$.
Writing $J_1^{\beta_j}=(\beta_j+2)t-(\beta_j+1)$ adds at most one
such factor. Multiplication by the fixed polynomial $a$ has bounded
angular shifts and bounded radial powers.
Finally
$2\nabla a\cdot\nabla U=\Delta(aU)-a\Delta U-(\Delta a)U$.
In the displayed Euler formula the large coefficient $p+j$ occurs
on the harmonic layer, whose Laplacian product costs only one
additional factor $p+j$; the layer-one coefficient is only $2j+2$.
The product therefore costs at most $C_a(p+j)^2$.
The direct Euler bound also gives $Cp^2$ after two fixed moments.

At the fixed finite chart, $q_0-1,d_0-1$ have coefficient multiplier
norm $O_N(\tau p^{-2})$. Put
\[
 \epsilon_{\mathrm{rec}}
 =\max\{\nrm{M_{q_0-1}}_{\Cc\to\Cc},
                 \nrm{M_{d_0-1}}_{\Cc\to\Cc}\},
\]
and increase the threshold until $\epsilon_{\mathrm{rec}}<1/2$.
Their inverse powers preserve all fixed moments: multiplication
by a polynomial of degree $D_N$ increases
degree by at most $D_N$, and
$\sum_{v\ge0}(1+vD_N)^k\epsilon_{\mathrm{rec}}^v<\infty$ for
$\epsilon_{\mathrm{rec}}<1/2$.
The constants may depend on $N,k$ but not on $p,r,\theta$.
Apply the two-layer estimates to every term of
\eqref{inf:scalar-pullback} and then these inverse multipliers.
Equation~\eqref{inf:lift-estimates} and $R^2=O(p^2)$ give the
first estimate in \eqref{inf:residual-moments}.
Apply \eqref{inf:Euler-moment}, using that first estimate at moment
$k+2$, to obtain the second.
The module bound for $\HP\eta$ and the fixed inverse multiplier
$d_0^{-1}$ then prove \eqref{inf:residual-material-bound}.
\end{proof}

\subsection{The exact material map and its inverse}

The radial harmonic chart specifies the interior extension as well
as the boundary displacement.

\begin{theorem}[Material isomorphism]\label{inf:material-isomorphism}
Fix $N\ge4$ and take $p$ sufficiently large for the centre above.
Put
\begin{equation}\label{inf:material-definitions}
 D_p=-\partial_{\nu\nu}U_0|_{\partial B},\qquad
 \mathfrak f_0=\frac{\Euler U_0}{d_0},\qquad
 T_\eta=(\HP\eta)\mathfrak f_0.
\end{equation}
The map
\begin{equation}\label{inf:Qmat-definition}
 \Qmat(\delta g,\eta)=\Delta(K\delta g-T_\eta)
                 :\Xc_p\longrightarrow\Cc
\end{equation}
is a bounded isomorphism, with
\begin{equation}\label{inf:Qmat-bounds}
 \nrm{\Qmat}\le C_Np^8,\qquad
 \nrm{\Qmat^{-1}}\le C_Np^6.
\end{equation}
If $J_D=B_{\Phi_N,D}\KD$, then the exact factorization is
\begin{equation}\label{inf:material-factorization}
 D\Fnl(z_0)=J_D\Qmat+\mathcal V_{\mathrm{res}}.
\end{equation}
All constants are uniform in the admissible split and
$0<|\theta|\le20$; the exponents are independent of $N$.
\end{theorem}
\begin{proof}
The centre has constant Dirichlet trace and zero gradient, so its
boundary Hessian is purely normal-normal. Evaluating its equation
at the boundary gives
\begin{equation}\label{inf:normal-hessian}
 D_p=\frac{R^2-E_0|_{\partial B}}{\nu^TA_{q_0}\nu}.
\end{equation}
Indeed all drift terms multiply the zero gradient, and
$D^2U_0=(\partial_{\nu\nu}U_0)\nu\otimes\nu$ there.
The finite centre polynomial and the boundary algebra show
$d_0|_{\partial B}=1+O_N(\tau p^{-2})$ and
$\nu^TA_{q_0}\nu=1+O_N(\tau p^{-2})$ in $\Ec_1$.
For the latter assertion insert \eqref{inf:pullback-metric};
the product formula in the proof of
\eqref{inf:gradient-multiplier}, with the extra boundary weight
$1+k\le(1+i)(1+j)$ on each angular product, gives
\[
 \nrm{M_{q_0}|_{\partial B}}_{\Ec_1}
 \le Cp\left(\sum_j(1+j)^2 2^j\wt(2j)|f_{N,j}|\right)^2
 \le C_N\tau^2p^{-3}.
\]
Here all fixed moments of the finite centre shape are
$O_N(\tau p^{-2})$ by \eqref{inf:centre-profile}.
The trace of $E_0$ in $\Ec_1$ is bounded by its first source
moment, since $J_s^{\beta_j}(1)=1$.
To identify this coefficient trace, take polynomial sums $F_m$ of
$E_0$ and put $h_m=F_m|_{\partial B}$. The first source moment
gives $h_m\to h$ in $\Ec_1$, while the fourth source moment in
Proposition~\ref{inf:residual-estimates} and
\eqref{inf:Delta-moment} gives $\Delta F_m\to\Delta E_0$ in $L^2$.
Since $F_m-\HP h_m$ has zero Dirichlet trace,
\[
 F_m-\HP h_m=\KD\Delta F_m.
\]
At fixed $p$, harmonic extension and Dirichlet Poisson regularity
make both terms on the right of
$F_m=\HP h_m+\KD\Delta F_m$ converge in $H^2$.
Their $L^2$ limit is $E_0$, hence
$E_0=\HP h+\KD\Delta E_0\in H^2$ and its Sobolev trace is $h$.
Thus \eqref{inf:residual-moments} and the boundary algebra give
\begin{equation}\label{inf:normal-hessian-inverse}
 \nrm{D_p^{-1}}_{\Ec_1}\le C_Np^{-2}
\end{equation}
by the convergent inverse series in \eqref{inf:normal-hessian}.
In particular $D_p$ has no zero.

The moment bounds for the fixed centre and its reciprocal $d_0$
give
$\nrm{\mathfrak f_0}_{\Cc^{[k]}}\le C_{N,k}p^{k+4}$ by
\eqref{inf:Euler-moment}.
Using \eqref{inf:Delta-moment} with its sum of two powers yields
\begin{equation}\label{inf:fixed-material-field}
 \nrm{\Delta \mathfrak f_0}_\Cc\le C_Np^8.
\end{equation}
Since $\mathfrak f_0$ has zero Dirichlet trace, Dirichlet uniqueness implies
$\mathfrak f_0=\KD\Delta \mathfrak f_0$.
Harmonicity of $\HP\eta$ and
\eqref{inf:gradient-C-one} therefore give
\begin{equation}\label{inf:material-field-bound}
 \nrm{\Delta T_\eta}_\Cc
 \le\nrm{M_{\HP\eta}\Delta \mathfrak f_0}_\Cc+
       2\nrm{\nabla(\HP\eta)\cdot\nabla\KD\Delta \mathfrak f_0}_\Cc
 \le C_Np^8\nrm\eta_{\Ec_1}.
\end{equation}
One moment of the unknown shape enters. This proves
boundedness of \eqref{inf:Qmat-definition}.

For an arbitrary source $\psi\in\Cc$ set
\begin{equation}\label{inf:material-reconstruction}
 W=\KD\psi,\quad N_W=\partial_\nu W,\quad
 \eta=(d_0|_{\partial B})N_W/D_p,\quad
 \delta g=\psi+\Delta T_\eta.
\end{equation}
Equations \eqref{inf:poisson-trace} and
\eqref{inf:normal-hessian-inverse} imply
$\nrm\eta_{\Ec_1}\le C_Np^{-2}\nrm\psi_\Cc$;
then \eqref{inf:material-field-bound} gives
$\nrm{\delta g}_\Cc\le(1+C_Np^6)\nrm\psi_\Cc$.
At the boundary,
\[
 T_\eta=0,\qquad \partial_\nu T_\eta=-D_p\eta/d_0.
\]
This follows by differentiating $\mathfrak f_0=(\Euler U_0)/d_0$:
$\Euler U_0=0$ there and
$\partial_\nu\Euler U_0=\partial_{\nu\nu}U_0=-D_p$.
Hence $W+T_\eta$ has both traces zero.
It is the Dirichlet Poisson solution with source $\delta g$;
\eqref{inf:poisson-trace} implies $\delta g\in\Gc_p$ and
$K\delta g=W+T_\eta$.
This proves that \eqref{inf:material-reconstruction} is a right
inverse to $\Qmat$ and gives its $p^6$ bound.
Conversely, for any input to $\Qmat$, the field
$W=K\delta g-T_\eta$ has zero Dirichlet trace and is recovered
uniquely from its Laplacian. Its normal trace is $D_p\eta/d_0$,
so \eqref{inf:material-reconstruction} recovers $\eta$ and then
$\delta g$. This proves the left inverse as well.

It remains to prove \eqref{inf:material-factorization}.
For the chart variation $\Phi_N+t(\HP\eta)y$, its reference material
vector is
$v_\eta=D\Phi_N^{-1}((\HP\eta)y)=y(\HP\eta)/d_0$.
Differentiation of the two compositions in
\eqref{inf:scalar-pullback-definition} gives the exact scalar
commutator identity
\[
 \delta(P_\Phi\Delta)U_0
 =-P_{\Phi_N}\Delta(v_\eta\cdot\nabla U_0)
       +v_\eta\cdot\nabla(P_{\Phi_N}\Delta U_0).
\]
It follows either by differentiating the inverse chart or by
applying the chain rule first to a smooth physical test function.
Since $v_\eta\cdot\nabla U_0=T_\eta$, adding the frequency term
and the source variation yields
\[
 D\Fnl(z_0)(\delta g,\eta)
 =B_{\Phi_N}(K\delta g-T_\eta)
       +(\HP\eta)d_0^{-1}\Euler E_0.
\]
The field in parentheses has zero Dirichlet trace, so it equals
$\KD\Qmat(\delta g,\eta)$. This proves
\eqref{inf:material-factorization} on finite variations.
To pass to the completion, fix $p$ and the smooth real centre chart.
Dirichlet regularity on the ball gives a bounded map
$\KD:L^2(B)\to H^2(B)\cap H^1_0(B)$; the smooth coefficients in
\eqref{inf:scalar-pullback} give a bounded map
$B_{\Phi_N,D}:H^2(B)\cap H^1_0(B)\to L^2(B)$.
These continuity bounds use only ball regularity and the fixed
chart, independently of the Dirichlet inverse estimates below.
For $\dot z\in\Xc_p$, take finite variations
$\dot z_m\to\dot z$ in $\Xc_p$.
The already proved boundedness of $\Qmat$ gives
$\Qmat\dot z_m\to\Qmat\dot z$ in $\Cc$, hence in $L^2(B)$ by
Lemma~\ref{inf:space-embedding}. Consequently
$J_D\Qmat\dot z_m\to B_{\Phi_N,D}\KD\Qmat\dot z$ in $L^2(B)$.
The finite-variation identity also says
$J_D\Qmat\dot z_m=(D\Fnl(z_0)-\mathcal V_{\mathrm{res}})\dot z_m$.
Its right side converges in $\Cc$ to
$(D\Fnl(z_0)-\mathcal V_{\mathrm{res}})\dot z$, by
Proposition~\ref{inf:nonlinear-map} and
\eqref{inf:residual-material-bound}. The continuous injective
embedding $\Cc\hookrightarrow L^2(B)$ identifies the two limits.
Thus $J_D\Qmat\dot z$ belongs to $\Cc$ and satisfies
\eqref{inf:material-factorization} for every $\dot z\in\Xc_p$.
Surjectivity of $\Qmat$ also gives boundedness of $J_D$ on all
of $\Cc$, with
$J_D=(D\Fnl(z_0)-\mathcal V_{\mathrm{res}})\Qmat^{-1}$.
\end{proof}
